%% AAAI 2027 Reproducibility Checklist -- filled for:
%% "Beyond a Single Verdict: Multi-Dimensional Bias Fingerprints for Responsible VLM Deployment"
%% Compile standalone (pdflatex) or %% AAAI 2027 Reproducibility Checklist -- filled for:
%% "Beyond a Single Verdict: Multi-Dimensional Bias Fingerprints for Responsible VLM Deployment"
%% Compile standalone (pdflatex) or %% AAAI 2027 Reproducibility Checklist -- filled for:
%% "Beyond a Single Verdict: Multi-Dimensional Bias Fingerprints for Responsible VLM Deployment"
%% Compile standalone (pdflatex) or %% AAAI 2027 Reproducibility Checklist -- filled for:
%% "Beyond a Single Verdict: Multi-Dimensional Bias Fingerprints for Responsible VLM Deployment"
%% Compile standalone (pdflatex) or \input{checklist} at the end of main.tex.
\documentclass[letterpaper]{article}
\usepackage[submission]{aaai25}
\usepackage{times}
\usepackage{helvet}
\usepackage{courier}
\usepackage[hyphens]{url}
\usepackage{graphicx}
\urlstyle{rm}
\def\UrlFont{\rm}
\usepackage{natbib}
\usepackage{caption}
\usepackage{enumitem}
\frenchspacing
\setlength{\pdfpagewidth}{8.5in}
\setlength{\pdfpageheight}{11in}
\pdfinfo{/TemplateVersion (2025.1)}

\title{Reproducibility Checklist\\ \large Beyond a Single Verdict: Multi-Dimensional Bias Fingerprints for Responsible VLM Deployment}
\author{}
\affiliations{}

\begin{document}
\maketitle

\noindent This checklist follows the AAAI Reproducibility Checklist. Answers
refer to the accompanying paper; section names below point to where each item is
substantiated.

\section*{1. General Paper Structure}
\begin{enumerate}[leftmargin=*]
\item Includes a conceptual outline and/or pseudocode description of AI methods introduced. \textbf{yes} --- the Deterministic Scoring Pipeline section gives the six scoring functions in closed form; the appendix provides the verbatim reference implementation of \texttt{score\_response} and all lexicons.
\item Clearly delineates statements that are opinions, hypothesis, and speculation from objective facts and results. \textbf{yes} --- the Motivation section frames the five premises as arguments; empirical claims are confined to Results and the appendix tables.
\item Provides well-marked pedagogical references for less-familiar readers to gain background necessary to replicate the paper. \textbf{yes} --- Related Work and the probe grounding cite the Stereotype Content Model and implicit-association literature.
\end{enumerate}

\section*{2. Theoretical Contributions}
\noindent \textbf{Does this paper make theoretical contributions? no.} The paper's
contribution is a benchmark and an empirical audit, together with a
measurement-methodology argument; it proves no formal theorems. Items 2a--2h are
therefore not applicable.

\section*{3. Datasets}
\noindent \textbf{Does this paper rely on one or more datasets? yes.}
\begin{enumerate}[leftmargin=*]
\item A motivation is given for why the experiments are conducted on the selected datasets. \textbf{yes} --- the Dataset subsection explains that FHIBE provides consented, self-reported demographic attributes, avoiding the circularity of imputed labels.
\item All novel datasets introduced in this paper are included in a data appendix. \textbf{NA} --- no novel dataset is introduced; the paper evaluates on the existing FHIBE corpus.
\item All novel datasets will be made publicly available upon publication with a license that allows free use for research. \textbf{NA} --- no novel dataset introduced.
\item All datasets drawn from the existing literature are accompanied by appropriate citation and, when available, a link/DOI. \textbf{yes} --- FHIBE is cited; the corpus statistics appendix documents counts and provenance.
\item All datasets that are not publicly available are described in detail, with an explanation of why publicly available alternatives are not used. \textbf{NA} --- FHIBE is the sole corpus and is a released, access-controlled human-centric benchmark; its access terms are described in the Ethics section.
\end{enumerate}

\section*{4. Computational Experiments}
\noindent \textbf{Does this paper include computational experiments? yes.}
\begin{enumerate}[leftmargin=*]
\item This paper states the number and range of values tried per (hyper-)parameter during development, along with the criterion used for selecting the final value. \textbf{yes} --- decoding uses the documented defaults (temperature 0.7, top-$p$ 0.9, max 256 tokens); scoring has no learned hyperparameters (deterministic lexicon rules), stated in Experimental Setup and the appendix.
\item Any code required for pre-processing data is included in the appendix. \textbf{yes} --- the appendix gives the verbatim scoring reference implementation and preprocessing description (seed 42).
\item All source code required for conducting and analyzing the experiments is included in a code appendix. \textbf{partial} --- the scoring/analysis reference code and all lexicons are in the appendix; the full inference-harness repository is not bundled in the anonymous submission but the deterministic pipeline is fully specified.
\item All source code required for conducting and analyzing the experiments will be made publicly available upon publication. \textbf{yes} --- the benchmark code and analysis scripts are intended for public release on acceptance.
\item All source code implementing new methods have comments detailing the implementation, with references to the paper where each step comes from. \textbf{yes} --- the reference implementation in the appendix is annotated and cross-referenced to the scoring definitions.
\item If an algorithm depends on randomness, then the method used for setting seeds is described in a way sufficient to allow replication of results. \textbf{yes} --- bootstrap resampling uses a fixed seed (42), stated in the Disparity Metrics section and the bootstrap-CI appendix; the scoring pipeline itself is deterministic (bit-exact across runs).
\item This paper specifies the computing infrastructure used for running experiments (hardware and software), including GPU/CPU models; amount of memory; operating system; names and versions of relevant software libraries and frameworks. \textbf{partial} --- Experimental Setup states NVIDIA A100 GPUs and 4-bit quantised inference; exact driver/library versions are not enumerated in the anonymous submission.
\item This paper formally describes evaluation metrics used and explains the motivation for choosing these metrics. \textbf{yes} --- per-probe max--min disparity, the composite, Cohen's $d$, and the Kruskal--Wallis $H$ statistic are defined formally with motivation in the Disparity Metrics section.
\item This paper states the number of algorithm runs used to compute each reported result. \textbf{yes} --- the full corpus (35{,}189 images $\times$ 5 probes = 175{,}945 image--probe pairs per model, 527{,}835 scored responses across three models) is evaluated once per model; scoring is deterministic, so repeated runs are identical.
\item Analysis of experiments goes beyond single-dimensional summaries of performance (e.g., average; median) to include measures of variation, confidence, or other distributional information. \textbf{yes} --- the appendix reports Kruskal--Wallis omnibus tests, Dunn's post-hoc contrasts with Holm correction, rank effect sizes ($\eta^2_H$), and bootstrap 95\% confidence intervals.
\item The significance of any improvement or decrease in performance is judged using appropriate statistical tests (e.g., Wilcoxon signed-rank). \textbf{yes} --- Kruskal--Wallis with Bonferroni correction across 15 model--probe pairs ($\alpha_{\mathrm{adj}}\approx0.00067$), plus Dunn's post-hoc; the paper explicitly cautions that large-$n$ significance be read against effect sizes.
\item This paper lists all final (hyper-)parameters used for each model/algorithm in the paper's experiments. \textbf{yes} --- decoding parameters and the complete scoring specification (lexicons, refusal patterns, thresholds) are listed in the appendix.
\end{enumerate}

\section*{Honest caveats (also stated in the paper)}
\begin{itemize}[leftmargin=*]
\item The bootstrap resampling unit is the \emph{image}, not the subject, because
the result databases expose image identifiers but not subject identifiers; the
intervals are therefore image-level and may modestly understate within-subject
correlation (stated in the bootstrap-CI appendix).
\item \texttt{IDEFICS2-8B}'s disparities are computed on the 33.8\% of inputs that
returned scorable output (59{,}495 / 175{,}945) and are reported with an explicit
coverage caveat throughout; coverage itself is analysed as a fairness quantity in
the Coverage section.
\end{itemize}

\end{document}
 at the end of main.tex.
\documentclass[letterpaper]{article}
\usepackage[submission]{aaai25}
\usepackage{times}
\usepackage{helvet}
\usepackage{courier}
\usepackage[hyphens]{url}
\usepackage{graphicx}
\urlstyle{rm}
\def\UrlFont{\rm}
\usepackage{natbib}
\usepackage{caption}
\usepackage{enumitem}
\frenchspacing
\setlength{\pdfpagewidth}{8.5in}
\setlength{\pdfpageheight}{11in}
\pdfinfo{/TemplateVersion (2025.1)}

\title{Reproducibility Checklist\\ \large Beyond a Single Verdict: Multi-Dimensional Bias Fingerprints for Responsible VLM Deployment}
\author{}
\affiliations{}

\begin{document}
\maketitle

\noindent This checklist follows the AAAI Reproducibility Checklist. Answers
refer to the accompanying paper; section names below point to where each item is
substantiated.

\section*{1. General Paper Structure}
\begin{enumerate}[leftmargin=*]
\item Includes a conceptual outline and/or pseudocode description of AI methods introduced. \textbf{yes} --- the Deterministic Scoring Pipeline section gives the six scoring functions in closed form; the appendix provides the verbatim reference implementation of \texttt{score\_response} and all lexicons.
\item Clearly delineates statements that are opinions, hypothesis, and speculation from objective facts and results. \textbf{yes} --- the Motivation section frames the five premises as arguments; empirical claims are confined to Results and the appendix tables.
\item Provides well-marked pedagogical references for less-familiar readers to gain background necessary to replicate the paper. \textbf{yes} --- Related Work and the probe grounding cite the Stereotype Content Model and implicit-association literature.
\end{enumerate}

\section*{2. Theoretical Contributions}
\noindent \textbf{Does this paper make theoretical contributions? no.} The paper's
contribution is a benchmark and an empirical audit, together with a
measurement-methodology argument; it proves no formal theorems. Items 2a--2h are
therefore not applicable.

\section*{3. Datasets}
\noindent \textbf{Does this paper rely on one or more datasets? yes.}
\begin{enumerate}[leftmargin=*]
\item A motivation is given for why the experiments are conducted on the selected datasets. \textbf{yes} --- the Dataset subsection explains that FHIBE provides consented, self-reported demographic attributes, avoiding the circularity of imputed labels.
\item All novel datasets introduced in this paper are included in a data appendix. \textbf{NA} --- no novel dataset is introduced; the paper evaluates on the existing FHIBE corpus.
\item All novel datasets will be made publicly available upon publication with a license that allows free use for research. \textbf{NA} --- no novel dataset introduced.
\item All datasets drawn from the existing literature are accompanied by appropriate citation and, when available, a link/DOI. \textbf{yes} --- FHIBE is cited; the corpus statistics appendix documents counts and provenance.
\item All datasets that are not publicly available are described in detail, with an explanation of why publicly available alternatives are not used. \textbf{NA} --- FHIBE is the sole corpus and is a released, access-controlled human-centric benchmark; its access terms are described in the Ethics section.
\end{enumerate}

\section*{4. Computational Experiments}
\noindent \textbf{Does this paper include computational experiments? yes.}
\begin{enumerate}[leftmargin=*]
\item This paper states the number and range of values tried per (hyper-)parameter during development, along with the criterion used for selecting the final value. \textbf{yes} --- decoding uses the documented defaults (temperature 0.7, top-$p$ 0.9, max 256 tokens); scoring has no learned hyperparameters (deterministic lexicon rules), stated in Experimental Setup and the appendix.
\item Any code required for pre-processing data is included in the appendix. \textbf{yes} --- the appendix gives the verbatim scoring reference implementation and preprocessing description (seed 42).
\item All source code required for conducting and analyzing the experiments is included in a code appendix. \textbf{partial} --- the scoring/analysis reference code and all lexicons are in the appendix; the full inference-harness repository is not bundled in the anonymous submission but the deterministic pipeline is fully specified.
\item All source code required for conducting and analyzing the experiments will be made publicly available upon publication. \textbf{yes} --- the benchmark code and analysis scripts are intended for public release on acceptance.
\item All source code implementing new methods have comments detailing the implementation, with references to the paper where each step comes from. \textbf{yes} --- the reference implementation in the appendix is annotated and cross-referenced to the scoring definitions.
\item If an algorithm depends on randomness, then the method used for setting seeds is described in a way sufficient to allow replication of results. \textbf{yes} --- bootstrap resampling uses a fixed seed (42), stated in the Disparity Metrics section and the bootstrap-CI appendix; the scoring pipeline itself is deterministic (bit-exact across runs).
\item This paper specifies the computing infrastructure used for running experiments (hardware and software), including GPU/CPU models; amount of memory; operating system; names and versions of relevant software libraries and frameworks. \textbf{partial} --- Experimental Setup states NVIDIA A100 GPUs and 4-bit quantised inference; exact driver/library versions are not enumerated in the anonymous submission.
\item This paper formally describes evaluation metrics used and explains the motivation for choosing these metrics. \textbf{yes} --- per-probe max--min disparity, the composite, Cohen's $d$, and the Kruskal--Wallis $H$ statistic are defined formally with motivation in the Disparity Metrics section.
\item This paper states the number of algorithm runs used to compute each reported result. \textbf{yes} --- the full corpus (35{,}189 images $\times$ 5 probes = 175{,}945 image--probe pairs per model, 527{,}835 scored responses across three models) is evaluated once per model; scoring is deterministic, so repeated runs are identical.
\item Analysis of experiments goes beyond single-dimensional summaries of performance (e.g., average; median) to include measures of variation, confidence, or other distributional information. \textbf{yes} --- the appendix reports Kruskal--Wallis omnibus tests, Dunn's post-hoc contrasts with Holm correction, rank effect sizes ($\eta^2_H$), and bootstrap 95\% confidence intervals.
\item The significance of any improvement or decrease in performance is judged using appropriate statistical tests (e.g., Wilcoxon signed-rank). \textbf{yes} --- Kruskal--Wallis with Bonferroni correction across 15 model--probe pairs ($\alpha_{\mathrm{adj}}\approx0.00067$), plus Dunn's post-hoc; the paper explicitly cautions that large-$n$ significance be read against effect sizes.
\item This paper lists all final (hyper-)parameters used for each model/algorithm in the paper's experiments. \textbf{yes} --- decoding parameters and the complete scoring specification (lexicons, refusal patterns, thresholds) are listed in the appendix.
\end{enumerate}

\section*{Honest caveats (also stated in the paper)}
\begin{itemize}[leftmargin=*]
\item The bootstrap resampling unit is the \emph{image}, not the subject, because
the result databases expose image identifiers but not subject identifiers; the
intervals are therefore image-level and may modestly understate within-subject
correlation (stated in the bootstrap-CI appendix).
\item \texttt{IDEFICS2-8B}'s disparities are computed on the 33.8\% of inputs that
returned scorable output (59{,}495 / 175{,}945) and are reported with an explicit
coverage caveat throughout; coverage itself is analysed as a fairness quantity in
the Coverage section.
\end{itemize}

\end{document}
 at the end of main.tex.
\documentclass[letterpaper]{article}
\usepackage[submission]{aaai25}
\usepackage{times}
\usepackage{helvet}
\usepackage{courier}
\usepackage[hyphens]{url}
\usepackage{graphicx}
\urlstyle{rm}
\def\UrlFont{\rm}
\usepackage{natbib}
\usepackage{caption}
\usepackage{enumitem}
\frenchspacing
\setlength{\pdfpagewidth}{8.5in}
\setlength{\pdfpageheight}{11in}
\pdfinfo{/TemplateVersion (2025.1)}

\title{Reproducibility Checklist\\ \large Beyond a Single Verdict: Multi-Dimensional Bias Fingerprints for Responsible VLM Deployment}
\author{}
\affiliations{}

\begin{document}
\maketitle

\noindent This checklist follows the AAAI Reproducibility Checklist. Answers
refer to the accompanying paper; section names below point to where each item is
substantiated.

\section*{1. General Paper Structure}
\begin{enumerate}[leftmargin=*]
\item Includes a conceptual outline and/or pseudocode description of AI methods introduced. \textbf{yes} --- the Deterministic Scoring Pipeline section gives the six scoring functions in closed form; the appendix provides the verbatim reference implementation of \texttt{score\_response} and all lexicons.
\item Clearly delineates statements that are opinions, hypothesis, and speculation from objective facts and results. \textbf{yes} --- the Motivation section frames the five premises as arguments; empirical claims are confined to Results and the appendix tables.
\item Provides well-marked pedagogical references for less-familiar readers to gain background necessary to replicate the paper. \textbf{yes} --- Related Work and the probe grounding cite the Stereotype Content Model and implicit-association literature.
\end{enumerate}

\section*{2. Theoretical Contributions}
\noindent \textbf{Does this paper make theoretical contributions? no.} The paper's
contribution is a benchmark and an empirical audit, together with a
measurement-methodology argument; it proves no formal theorems. Items 2a--2h are
therefore not applicable.

\section*{3. Datasets}
\noindent \textbf{Does this paper rely on one or more datasets? yes.}
\begin{enumerate}[leftmargin=*]
\item A motivation is given for why the experiments are conducted on the selected datasets. \textbf{yes} --- the Dataset subsection explains that FHIBE provides consented, self-reported demographic attributes, avoiding the circularity of imputed labels.
\item All novel datasets introduced in this paper are included in a data appendix. \textbf{NA} --- no novel dataset is introduced; the paper evaluates on the existing FHIBE corpus.
\item All novel datasets will be made publicly available upon publication with a license that allows free use for research. \textbf{NA} --- no novel dataset introduced.
\item All datasets drawn from the existing literature are accompanied by appropriate citation and, when available, a link/DOI. \textbf{yes} --- FHIBE is cited; the corpus statistics appendix documents counts and provenance.
\item All datasets that are not publicly available are described in detail, with an explanation of why publicly available alternatives are not used. \textbf{NA} --- FHIBE is the sole corpus and is a released, access-controlled human-centric benchmark; its access terms are described in the Ethics section.
\end{enumerate}

\section*{4. Computational Experiments}
\noindent \textbf{Does this paper include computational experiments? yes.}
\begin{enumerate}[leftmargin=*]
\item This paper states the number and range of values tried per (hyper-)parameter during development, along with the criterion used for selecting the final value. \textbf{yes} --- decoding uses the documented defaults (temperature 0.7, top-$p$ 0.9, max 256 tokens); scoring has no learned hyperparameters (deterministic lexicon rules), stated in Experimental Setup and the appendix.
\item Any code required for pre-processing data is included in the appendix. \textbf{yes} --- the appendix gives the verbatim scoring reference implementation and preprocessing description (seed 42).
\item All source code required for conducting and analyzing the experiments is included in a code appendix. \textbf{partial} --- the scoring/analysis reference code and all lexicons are in the appendix; the full inference-harness repository is not bundled in the anonymous submission but the deterministic pipeline is fully specified.
\item All source code required for conducting and analyzing the experiments will be made publicly available upon publication. \textbf{yes} --- the benchmark code and analysis scripts are intended for public release on acceptance.
\item All source code implementing new methods have comments detailing the implementation, with references to the paper where each step comes from. \textbf{yes} --- the reference implementation in the appendix is annotated and cross-referenced to the scoring definitions.
\item If an algorithm depends on randomness, then the method used for setting seeds is described in a way sufficient to allow replication of results. \textbf{yes} --- bootstrap resampling uses a fixed seed (42), stated in the Disparity Metrics section and the bootstrap-CI appendix; the scoring pipeline itself is deterministic (bit-exact across runs).
\item This paper specifies the computing infrastructure used for running experiments (hardware and software), including GPU/CPU models; amount of memory; operating system; names and versions of relevant software libraries and frameworks. \textbf{partial} --- Experimental Setup states NVIDIA A100 GPUs and 4-bit quantised inference; exact driver/library versions are not enumerated in the anonymous submission.
\item This paper formally describes evaluation metrics used and explains the motivation for choosing these metrics. \textbf{yes} --- per-probe max--min disparity, the composite, Cohen's $d$, and the Kruskal--Wallis $H$ statistic are defined formally with motivation in the Disparity Metrics section.
\item This paper states the number of algorithm runs used to compute each reported result. \textbf{yes} --- the full corpus (35{,}189 images $\times$ 5 probes = 175{,}945 image--probe pairs per model, 527{,}835 scored responses across three models) is evaluated once per model; scoring is deterministic, so repeated runs are identical.
\item Analysis of experiments goes beyond single-dimensional summaries of performance (e.g., average; median) to include measures of variation, confidence, or other distributional information. \textbf{yes} --- the appendix reports Kruskal--Wallis omnibus tests, Dunn's post-hoc contrasts with Holm correction, rank effect sizes ($\eta^2_H$), and bootstrap 95\% confidence intervals.
\item The significance of any improvement or decrease in performance is judged using appropriate statistical tests (e.g., Wilcoxon signed-rank). \textbf{yes} --- Kruskal--Wallis with Bonferroni correction across 15 model--probe pairs ($\alpha_{\mathrm{adj}}\approx0.00067$), plus Dunn's post-hoc; the paper explicitly cautions that large-$n$ significance be read against effect sizes.
\item This paper lists all final (hyper-)parameters used for each model/algorithm in the paper's experiments. \textbf{yes} --- decoding parameters and the complete scoring specification (lexicons, refusal patterns, thresholds) are listed in the appendix.
\end{enumerate}

\section*{Honest caveats (also stated in the paper)}
\begin{itemize}[leftmargin=*]
\item The bootstrap resampling unit is the \emph{image}, not the subject, because
the result databases expose image identifiers but not subject identifiers; the
intervals are therefore image-level and may modestly understate within-subject
correlation (stated in the bootstrap-CI appendix).
\item \texttt{IDEFICS2-8B}'s disparities are computed on the 33.8\% of inputs that
returned scorable output (59{,}495 / 175{,}945) and are reported with an explicit
coverage caveat throughout; coverage itself is analysed as a fairness quantity in
the Coverage section.
\end{itemize}

\end{document}
 at the end of main.tex.
\documentclass[letterpaper]{article}
\usepackage[submission]{aaai25}
\usepackage{times}
\usepackage{helvet}
\usepackage{courier}
\usepackage[hyphens]{url}
\usepackage{graphicx}
\urlstyle{rm}
\def\UrlFont{\rm}
\usepackage{natbib}
\usepackage{caption}
\usepackage{enumitem}
\frenchspacing
\setlength{\pdfpagewidth}{8.5in}
\setlength{\pdfpageheight}{11in}
\pdfinfo{/TemplateVersion (2025.1)}

\title{Reproducibility Checklist\\ \large Beyond a Single Verdict: Multi-Dimensional Bias Fingerprints for Responsible VLM Deployment}
\author{}
\affiliations{}

\begin{document}
\maketitle

\noindent This checklist follows the AAAI Reproducibility Checklist. Answers
refer to the accompanying paper; section names below point to where each item is
substantiated.

\section*{1. General Paper Structure}
\begin{enumerate}[leftmargin=*]
\item Includes a conceptual outline and/or pseudocode description of AI methods introduced. \textbf{yes} --- the Deterministic Scoring Pipeline section gives the six scoring functions in closed form; the appendix provides the verbatim reference implementation of \texttt{score\_response} and all lexicons.
\item Clearly delineates statements that are opinions, hypothesis, and speculation from objective facts and results. \textbf{yes} --- the Motivation section frames the five premises as arguments; empirical claims are confined to Results and the appendix tables.
\item Provides well-marked pedagogical references for less-familiar readers to gain background necessary to replicate the paper. \textbf{yes} --- Related Work and the probe grounding cite the Stereotype Content Model and implicit-association literature.
\end{enumerate}

\section*{2. Theoretical Contributions}
\noindent \textbf{Does this paper make theoretical contributions? no.} The paper's
contribution is a benchmark and an empirical audit, together with a
measurement-methodology argument; it proves no formal theorems. Items 2a--2h are
therefore not applicable.

\section*{3. Datasets}
\noindent \textbf{Does this paper rely on one or more datasets? yes.}
\begin{enumerate}[leftmargin=*]
\item A motivation is given for why the experiments are conducted on the selected datasets. \textbf{yes} --- the Dataset subsection explains that FHIBE provides consented, self-reported demographic attributes, avoiding the circularity of imputed labels.
\item All novel datasets introduced in this paper are included in a data appendix. \textbf{NA} --- no novel dataset is introduced; the paper evaluates on the existing FHIBE corpus.
\item All novel datasets will be made publicly available upon publication with a license that allows free use for research. \textbf{NA} --- no novel dataset introduced.
\item All datasets drawn from the existing literature are accompanied by appropriate citation and, when available, a link/DOI. \textbf{yes} --- FHIBE is cited; the corpus statistics appendix documents counts and provenance.
\item All datasets that are not publicly available are described in detail, with an explanation of why publicly available alternatives are not used. \textbf{NA} --- FHIBE is the sole corpus and is a released, access-controlled human-centric benchmark; its access terms are described in the Ethics section.
\end{enumerate}

\section*{4. Computational Experiments}
\noindent \textbf{Does this paper include computational experiments? yes.}
\begin{enumerate}[leftmargin=*]
\item This paper states the number and range of values tried per (hyper-)parameter during development, along with the criterion used for selecting the final value. \textbf{yes} --- decoding uses the documented defaults (temperature 0.7, top-$p$ 0.9, max 256 tokens); scoring has no learned hyperparameters (deterministic lexicon rules), stated in Experimental Setup and the appendix.
\item Any code required for pre-processing data is included in the appendix. \textbf{yes} --- the appendix gives the verbatim scoring reference implementation and preprocessing description (seed 42).
\item All source code required for conducting and analyzing the experiments is included in a code appendix. \textbf{partial} --- the scoring/analysis reference code and all lexicons are in the appendix; the full inference-harness repository is not bundled in the anonymous submission but the deterministic pipeline is fully specified.
\item All source code required for conducting and analyzing the experiments will be made publicly available upon publication. \textbf{yes} --- the benchmark code and analysis scripts are intended for public release on acceptance.
\item All source code implementing new methods have comments detailing the implementation, with references to the paper where each step comes from. \textbf{yes} --- the reference implementation in the appendix is annotated and cross-referenced to the scoring definitions.
\item If an algorithm depends on randomness, then the method used for setting seeds is described in a way sufficient to allow replication of results. \textbf{yes} --- bootstrap resampling uses a fixed seed (42), stated in the Disparity Metrics section and the bootstrap-CI appendix; the scoring pipeline itself is deterministic (bit-exact across runs).
\item This paper specifies the computing infrastructure used for running experiments (hardware and software), including GPU/CPU models; amount of memory; operating system; names and versions of relevant software libraries and frameworks. \textbf{partial} --- Experimental Setup states NVIDIA A100 GPUs and 4-bit quantised inference; exact driver/library versions are not enumerated in the anonymous submission.
\item This paper formally describes evaluation metrics used and explains the motivation for choosing these metrics. \textbf{yes} --- per-probe max--min disparity, the composite, Cohen's $d$, and the Kruskal--Wallis $H$ statistic are defined formally with motivation in the Disparity Metrics section.
\item This paper states the number of algorithm runs used to compute each reported result. \textbf{yes} --- the full corpus (35{,}189 images $\times$ 5 probes = 175{,}945 image--probe pairs per model, 527{,}835 scored responses across three models) is evaluated once per model; scoring is deterministic, so repeated runs are identical.
\item Analysis of experiments goes beyond single-dimensional summaries of performance (e.g., average; median) to include measures of variation, confidence, or other distributional information. \textbf{yes} --- the appendix reports Kruskal--Wallis omnibus tests, Dunn's post-hoc contrasts with Holm correction, rank effect sizes ($\eta^2_H$), and bootstrap 95\% confidence intervals.
\item The significance of any improvement or decrease in performance is judged using appropriate statistical tests (e.g., Wilcoxon signed-rank). \textbf{yes} --- Kruskal--Wallis with Bonferroni correction across 15 model--probe pairs ($\alpha_{\mathrm{adj}}\approx0.00067$), plus Dunn's post-hoc; the paper explicitly cautions that large-$n$ significance be read against effect sizes.
\item This paper lists all final (hyper-)parameters used for each model/algorithm in the paper's experiments. \textbf{yes} --- decoding parameters and the complete scoring specification (lexicons, refusal patterns, thresholds) are listed in the appendix.
\end{enumerate}

\section*{Honest caveats (also stated in the paper)}
\begin{itemize}[leftmargin=*]
\item The bootstrap resampling unit is the \emph{image}, not the subject, because
the result databases expose image identifiers but not subject identifiers; the
intervals are therefore image-level and may modestly understate within-subject
correlation (stated in the bootstrap-CI appendix).
\item \texttt{IDEFICS2-8B}'s disparities are computed on the 33.8\% of inputs that
returned scorable output (59{,}495 / 175{,}945) and are reported with an explicit
coverage caveat throughout; coverage itself is analysed as a fairness quantity in
the Coverage section.
\end{itemize}

\end{document}
